
This work demonstrated that simplicity bias creates a detectable signature in training dynamics: minority-group samples exhibit $8\times$ higher loss at 5\% of training epochs compared to majority samples. This signal enables 6.$6\times$ improvement over random baseline for minority detection without group labels.

Linear probe analysis established the mechanism: model representations encode spurious features with approximately 10 percentage points higher accuracy than core features throughout training. With pretrained features, this manifests as an accuracy magnitude gap rather than a temporal gap---both feature types peak at the same epoch.

The detection precision of 33\% reflects base rate constraints rather than signal weakness; the statistical separation is highly significant (p < $10^{-14})$. This suggests that loss-based detection could inform weighted training schemes, though intervention effectiveness remains to be validated.

